\section*{2. A complement to XVIII}

\subsubsection*{Proposition 2.1.}
Let $X\subset\mathbb P$ be a smooth projective variety, pure of odd
dimension $2n+1$ over an algebraically closed field $k$, and let
$\ell\subset\check{\mathbb P}$ be a generic Lefschetz pencil of hyperplane
sections of $X$.  Let $s$ be a geometric point of $\ell$, let $H_s$ be the
corresponding hyperplane section, let $S\subset\ell$ be the set of exceptional
values of the pencil, and let $\pi_1(\ell-S,s)$ be the monodromy group.
Suppose that $X$ satisfies condition (LV).  Then the representation of
$\pi_1(\ell-S,s)$ on the vanishing cohomology
\[
\operatorname{Ev}(H_s,\mathbb Q_\ell)
\subset H^{2n}(H_s,\mathbb Q_\ell)
\]
is absolutely irreducible.\footnote{The print says that $X$ has even
dimension $2n$.  The displayed degree $H^{2n}(H_s)$, the reflection formula,
Remark 2.2, and the application to Theorem 1.4 all force odd dimension
$2n+1$.}

Recall that $\operatorname{Ev}(H_s,\mathbb Q_\ell)$ is the orthogonal
complement of the image of $H^{2n}(X,\mathbb Q_\ell)$ in
$H^{2n}(H_s,\mathbb Q_\ell)$, and that this image equals the subgroup of
monodromy invariants (XVIII 5.6.10).

\subsubsection*{Proof.}
The projective line $\mathbb P^1$ is simply connected.  Hence
$\pi_1(\ell-S,s)$ is topologically generated by the inertia subgroups at the
points of $S$ and their conjugates.  Its image in
$\operatorname{GL}(H^{2n}(H_s,\mathbb Q_\ell))$ is therefore topologically
generated by the transformations
\[
x\longmapsto x-(-1)^n(x\cdot\delta)\delta,
\]
where $\delta$ is conjugate to a vanishing cycle.\footnote{The print gives
$H^n(H_s)$ at this point, although every surrounding occurrence and the
statement of the proposition use the middle degree $H^{2n}(H_s)$.}
Consequently, the image of $H^{2n}(X,\mathbb Q_\ell)$ in
$H^{2n}(H_s,\mathbb Q_\ell)$ is the orthogonal complement of the subspace
generated by the vanishing cycles and their conjugates; thus
$\operatorname{Ev}(H_s,\mathbb Q_\ell)$ is generated by those cycles and
their conjugates.

For a \underline{generic} Lefschetz pencil, the arguments of XVIII §6 show
that the vanishing cycles attached to the different points of $S$ are
conjugate to one another.

Let $W\subset\operatorname{Ev}(H_s,\mathbb Q_\ell)$ be a nonzero invariant
subspace.  There exist a conjugate $\delta$ of a vanishing cycle and an
element $v\in W$ such that $(v,\delta)\ne0$; this is where condition (LV) is
used.  Since
\[
v-(-1)^n(v\cdot\delta)\delta\in W,
\]
we have $\delta\in W$.\footnote{The print uses the undefined exponent $m$ in
this second occurrence; the immediately preceding reflection and the
Picard--Lefschetz formula give $n$.}  The subspace $W$, containing one
vanishing cycle and all its conjugates, is therefore
$\operatorname{Ev}(H_s,\mathbb Q_\ell)$.  The argument remains valid after
extending the coefficient field $\mathbb Q_\ell$, proving absolute
irreducibility.

\subsubsection*{Remark 2.2.}
If condition (LV) is omitted, the preceding arguments still show that every
invariant subspace of $H^{2n}(H_s,\mathbb Q_\ell)$ either contains
$\operatorname{Ev}(H_s,\mathbb Q_\ell)$ or is contained in the image of
$H^{2n}(X,\mathbb Q_\ell)$.  The same assertion holds for $X$ of even
dimension (and for every Lefschetz pencil).

